% concatenated from quartic_arxiv.tex
\documentclass[12pt]{article}

\usepackage[margin=1in]{geometry}
\usepackage{amsmath, amsthm, amssymb}
\usepackage{microtype}
\usepackage{hyperref}

% Sequential numbering: Theorem 1, Lemma 2, Proposition 3, etc.
\newtheorem{theorem}{Theorem}
\newtheorem{lemma}[theorem]{Lemma}
\newtheorem{proposition}[theorem]{Proposition}
\newtheorem{corollary}[theorem]{Corollary}
\theoremstyle{definition}
\newtheorem{definition}[theorem]{Definition}
\newtheorem{example}[theorem]{Example}
\theoremstyle{remark}
\newtheorem{remark}[theorem]{Remark}

\title{A Simple Trigonometric Classification of Quartic Roots}

\author{Sawon Pratiher\footnote{Indian Institute of Technology (IIT) Kharagpur, West Bengal 721302, India. This work is the outcome of the culmination of research started since the author was first introduced to the algebra coursework in school. (E-mail: \texttt{sawon1234@gmail.com})}}

\date{March 17, 2022}

\begin{document}

\maketitle

\begin{abstract}
This article provides a simple trigonometric method for determining how many roots of a quartic equation are real and how many are complex, without solving the equation. The approach replaces the quartic's classical discriminant---a degree-six polynomial in the coefficients---with an elementary analysis of the function $f(\theta) = a\cos\theta + \cos 4\theta + b$ on $[0,\pi]$, obtained by matching the quartic to the Chebyshev identity $8\cos^4\!\theta - 8\cos^2\!\theta + 1 = \cos 4\theta$. The derivation is computationally light and conceptually natural, and has the potential to demystify the geometry of a quartic equation's roots from a trigonometric perspective.
\end{abstract}

\section{Introduction}

Determining whether a polynomial equation has all real roots, or some complex ones, is a question as old as algebra itself. For the quadratic $x^2 + bx + c = 0$, the answer is immediate: compute $b^2 - 4c$ and check its sign. For the cubic, the discriminant is a bit more involved but still manageable. For the quartic, however, the discriminant is a polynomial of degree six in the coefficients~\cite{Rees1922}, and its evaluation is a formidable exercise even for those who have seen it before.

This article introduces a simple alternative. We show that a single trigonometric substitution---$t = u\cos\theta$ for a well-chosen constant $u$---transforms the quartic into a function of the form $a\cos\theta + \cos 4\theta + b$, and that counting the zeros of this function on $[0,\pi]$ immediately tells us how many roots of the quartic are real. The derivation uses nothing beyond the Chebyshev identity for $\cos 4\theta$, the intermediate value theorem, and the convexity of the quartic outside a bounded interval. Every step is a routine computation that a first-time Algebra learner could follow.

The reader may wonder: why should a trigonometric function know anything about the roots of a polynomial? The connection is the identity
\begin{equation}\label{eq:cheb4}
8\cos^4\theta - 8\cos^2\theta + 1 = \cos 4\theta,
\end{equation}
which is nothing other than the statement that the degree-four Chebyshev polynomial $T_4(x) = 8x^4 - 8x^2 + 1$ satisfies $T_4(\cos\theta) = \cos 4\theta$. If we substitute $t = u\cos\theta$ into a quartic with the right structure, the leading term $t^4$ and the quadratic term $t^2$ can be made to match the left side of (\ref{eq:cheb4}) exactly, leaving only $\cos\theta$ and a constant as residual terms. The bounded range of $\cos\theta$, together with the oscillatory nature of $\cos 4\theta$, then makes it possible to read off the root structure from an entirely elementary analysis.

Throughout this article, all polynomials have real coefficients, so complex roots come in conjugate pairs. A quartic with real coefficients therefore has zero, two, or four real roots.

\section{Derivation}

This section is intentionally written at a leisurely pace, to emphasize how straightforward all the algebraic manipulations are.

\subsection{The depressed quartic}

Every monic quartic $z^4 + a_3 z^3 + a_2 z^2 + a_1 z + a_0 = 0$ can be converted to a \emph{depressed} quartic---one with no cubic term---by the substitution $t = z - a_3/4$. This is a standard change of variables (see, e.g., \cite{Irving2013}), and the result is an equation of the form
\begin{equation}\label{eq:depressed}
P(t) = t^4 + mt^2 + pt + q = 0.
\end{equation}
By Vieta's relations~\cite{Viete1646}, the four roots $t_1, t_2, t_3, t_4$ of the depressed quartic satisfy $t_1 + t_2 + t_3 + t_4 = 0$, so the centroid of the four roots sits at the origin. In particular, knowing any three roots determines the fourth via $t_4 = -(t_1 + t_2 + t_3)$.

\subsection{The trigonometric substitution}

Here is the key step. Assume $m < 0$, and set $u = \sqrt{-m}$. (The case $m \geq 0$ is discussed in Section~\ref{sec:mgeq0}.) Substituting $t = u\cos\theta$ into (\ref{eq:depressed}) and dividing through by $u^4/8 = m^2/8$ gives
\begin{equation}\label{eq:after_sub}
8\cos^4\theta - 8\cos^2\theta + \frac{8p}{(-m)^{3/2}}\cos\theta + \frac{8q}{m^2} = 0.
\end{equation}
Observe that the first two terms are exactly the left side of the Chebyshev identity (\ref{eq:cheb4}). Subtracting that identity from (\ref{eq:after_sub}), we arrive at the \textbf{reduced trigonometric equation}:
\begin{equation}\label{eq:f_def}
\boxed{f(\theta) \;=\; a\cos\theta + \cos 4\theta + b \;=\; 0,}
\end{equation}
where the \textbf{trigonometric parameters} are
\begin{equation}\label{eq:ab}
a = \frac{8p}{(-m)^{3/2}}, \qquad b = \frac{8q}{m^2} - 1.
\end{equation}

That is the entire reduction. One substitution, one division, and one use of the Chebyshev identity. The quartic has been converted into a trigonometric equation in~$\theta$ with just two parameters, $a$ and $b$.

\subsection{Why this works: the bijection}

Since $\cos\colon [0,\pi] \to [-1,1]$ is a continuous, strictly decreasing bijection, the map $\theta \mapsto u\cos\theta$ is a bijection from $[0,\pi]$ onto $[-u, u]$. A direct calculation confirms the identity
\begin{equation}\label{eq:fP}
f(\theta) = \frac{8\,P(u\cos\theta)}{u^4}
\end{equation}
for all $\theta$, so the zeros of $f$ on $[0,\pi]$ correspond exactly to the roots of $P$ in $[-u, u]$. In particular, evaluating at the endpoints:
\begin{equation}\label{eq:boundary}
f(0) = \frac{8\,P(u)}{u^4}, \qquad f(\pi) = \frac{8\,P(-u)}{u^4}.
\end{equation}
The sign of $f(0)$ tells us whether the quartic is positive or negative at $t = u$, and the sign of $f(\pi)$ tells us the same at $t = -u$. In terms of the parameters $a$ and $b$:
\begin{equation}\label{eq:boundary_vals}
f(0) = a + 1 + b, \qquad f(\pi) = -a + 1 + b.
\end{equation}

\subsection{What about roots outside $[-u,u]$?}

The substitution $t = u\cos\theta$ only captures roots in $[-u,u]$, since $\cos\theta$ is bounded between $-1$ and $+1$. Could the quartic have real roots outside this interval? It certainly can. But a simple convexity argument controls how many.

\begin{lemma}[{Convexity outside $[-u,u]$}]\label{lem:convex}
If $m < 0$, then $P$ is strictly convex on $(-\infty, -u)$ and on $(u, \infty)$.
\end{lemma}

\begin{proof}
The second derivative is $P''(t) = 12t^2 + 2m$. For $|t| > u = \sqrt{-m}$, we have $t^2 > -m$, so $P''(t) > 12(-m) + 2m = -10m > 0$.
\end{proof}

A strictly convex function can cross the horizontal axis at most twice, but since $P(t) \to +\infty$ as $t \to \pm\infty$ for a quartic with positive leading coefficient, it actually crosses at most \emph{once} on each half-line. This yields:

\begin{corollary}\label{cor:ext}
The quartic $P$ has at most one root in $(u, \infty)$ and at most one root in $(-\infty, -u)$. Moreover:
\begin{itemize}
\item $P$ has exactly one root in $(u, \infty)$ if and only if $P(u) < 0$, i.e., $f(0) < 0$.
\item $P$ has exactly one root in $(-\infty, -u)$ if and only if $P(-u) < 0$, i.e., $f(\pi) < 0$.
\item $P$ has no roots outside $[-u,u]$ if and only if $f(0) \geq 0$ and $f(\pi) \geq 0$.
\end{itemize}
\end{corollary}

\begin{proof}
Since $P(t) \to +\infty$ as $t \to +\infty$ and $P$ is strictly convex and continuous on $(u, \infty)$, the intermediate value theorem guarantees exactly one crossing if $P(u) < 0$, and no crossing if $P(u) \geq 0$. The argument for $(-\infty, -u)$ is identical.
\end{proof}

So the problem neatly separates into two pieces: the \textbf{interior roots} (zeros of $f$ on $[0,\pi]$, corresponding to roots of $P$ in $[-u,u]$) and the \textbf{exterior roots} (at most one on each side, detected by the signs of $f(0)$ and $f(\pi)$).

\section{Classification of the roots}

We are now ready to state the main results. Define
\[
N_{\mathrm{int}} = \#\{\theta \in [0,\pi] : f(\theta) = 0\}, \qquad N_{\mathrm{ext}} = \mathbf{1}_{f(0)<0} + \mathbf{1}_{f(\pi)<0},
\]
where $\mathbf{1}$ denotes the indicator function. By Proposition~\ref{prop:bijection_formal} and Corollary~\ref{cor:ext}, the total number of real roots of $P$ is
\begin{equation}\label{eq:Nreal}
N_{\mathrm{real}} = N_{\mathrm{int}} + N_{\mathrm{ext}}.
\end{equation}

\begin{proposition}\label{prop:bijection_formal}
The zeros of $f$ in $[0,\pi]$ are in bijection with the roots of $P$ in $[-u,u]$. There is no double-counting at the boundary: if $f(0) = 0$ then $t = u$ is counted in $N_{\mathrm{int}}$, and $N_{\mathrm{ext}}$ contributes zero from the right (and similarly for $f(\pi) = 0$).
\end{proposition}

\subsection{Critical-point analysis}

To count the zeros of $f$ on $[0,\pi]$, we study its critical points. Differentiating (\ref{eq:f_def}):
\[
f'(\theta) = -a\sin\theta - 4\sin 4\theta.
\]
Using $\sin 4\theta = 4\sin\theta\cos\theta\cos 2\theta$ and noting that $\sin\theta > 0$ on $(0,\pi)$, the interior critical points satisfy
\begin{equation}\label{eq:crit_quartic}
a + 16\cos\theta\cos 2\theta = 0.
\end{equation}
Setting $x = \cos\theta \in (-1,1)$ and writing $\cos 2\theta = 2x^2 - 1$, this becomes
\begin{equation}\label{eq:h_cubic}
h(x) := x(2x^2 - 1) = -\frac{a}{16}.
\end{equation}

\begin{lemma}\label{lem:h_range}
The function $h(x) = 2x^3 - x$ on $[-1,1]$ has a local maximum of $\sqrt{6}/9$ at $x = -1/\sqrt{6}$ and a local minimum of $-\sqrt{6}/9$ at $x = 1/\sqrt{6}$. Its range on $[-1,1]$ is $[-1,1]$.
\end{lemma}

\begin{proof}
Setting $h'(x) = 6x^2 - 1 = 0$ gives $x = \pm 1/\sqrt{6}$. Direct evaluation yields $h(\pm 1/\sqrt{6}) = \mp\sqrt{6}/9$. At the endpoints, $h(\pm 1) = \pm 1$.
\end{proof}

It follows that $f$ has at most three interior critical points in $(0,\pi)$, so $f$ is piecewise monotone with at most four monotone segments, hence at most four zeros in $[0,\pi]$---consistent with the quartic having at most four real roots.

\subsection{The three cases}

The classification now follows from counting zeros of $f$ together with the boundary analysis.

\begin{theorem}[All four roots complex]\label{thm:four_complex}
Suppose $m < 0$. All four roots of $P(t) = 0$ are complex if and only if $f(\theta) > 0$ for all $\theta \in [0,\pi]$ (equivalently, $N_{\mathrm{int}} = 0$ and $N_{\mathrm{ext}} = 0$). A sufficient condition is
\begin{equation}\label{eq:four_complex_suff}
b > |a| + 1.
\end{equation}
\end{theorem}

\begin{proof}
If $f > 0$ on $[0,\pi]$, then $N_{\mathrm{int}} = 0$. The boundary values $f(0)$ and $f(\pi)$ are both positive, so $P(u) > 0$ and $P(-u) > 0$ by (\ref{eq:boundary}), giving $N_{\mathrm{ext}} = 0$. Thus $N_{\mathrm{real}} = 0$.

For the sufficient condition: for all $\theta$, $|a\cos\theta + \cos 4\theta| \leq |a| + 1$, so $f(\theta) \geq b - (|a|+1) > 0$ when $b > |a| + 1$.

Conversely, if $f < 0$ on $[0,\pi]$, then $N_{\mathrm{int}} = 0$ but $f(0) < 0$ and $f(\pi) < 0$, giving $N_{\mathrm{ext}} = 2$, so there are two (not zero) real roots. Thus the only route to zero real roots is $f > 0$.
\end{proof}

\begin{theorem}[Two real, two complex roots]\label{thm:two_real}
Suppose $m < 0$. The quartic $P(t) = 0$ has exactly two real roots if and only if $N_{\mathrm{int}} + N_{\mathrm{ext}} = 2$. This occurs generically in three scenarios:
\begin{itemize}
\item[(a)] $f < 0$ on $[0,\pi]$, so $N_{\mathrm{int}} = 0$ and $N_{\mathrm{ext}} = 2$ (one root on each side of $[-u,u]$). A sufficient condition is $b < -(|a|+1)$.
\item[(b)] $f$ has exactly two zeros in $[0,\pi]$ and $f(0) \geq 0$, $f(\pi) \geq 0$, so $N_{\mathrm{int}} = 2$ and $N_{\mathrm{ext}} = 0$.
\item[(c)] $f$ has exactly one zero in $[0,\pi]$ and exactly one of $f(0) < 0$ or $f(\pi) < 0$, so $N_{\mathrm{int}} = 1$ and $N_{\mathrm{ext}} = 1$.
\end{itemize}
\end{theorem}

\begin{theorem}[All four roots real]\label{thm:four_real}
Suppose $m < 0$. All four roots of $P(t) = 0$ are real if and only if $N_{\mathrm{int}} + N_{\mathrm{ext}} = 4$. When all four roots lie in $[-u,u]$ (equivalently, $f(0) \geq 0$ and $f(\pi) \geq 0$), they are given by
\[
t_k = u\cos\theta_k, \qquad k = 1, 2, 3, 4,
\]
where $0 \leq \theta_1 < \theta_2 < \theta_3 < \theta_4 \leq \pi$ are the four zeros of $f$ in $[0,\pi]$.
\end{theorem}

\section{Examples of use as a method}

The computational simplicity of the derivation renders it natural to use as a method: compute $a$ and $b$, examine $f(\theta) = a\cos\theta + \cos 4\theta + b$, and read off the root structure.

\begin{example}[Four real roots]\label{ex:four_real}
Consider $P(t) = t^4 - 25t^2 - 60t - 36 = 0$. Here $m = -25$, $p = -60$, $q = -36$, so $u = 5$. The trigonometric parameters are
\[
a = \frac{8(-60)}{125} = -3.84, \qquad b = \frac{8(-36)}{625} - 1 = -1.4608.
\]
The boundary values are $f(0) = -3.84 + 1 - 1.4608 = -4.3008 < 0$ and $f(\pi) = 3.84 + 1 - 1.4608 = 3.3792 > 0$. Since $f(0) < 0$, there is one root in $(5, \infty)$; since $f(\pi) > 0$, there are no roots in $(-\infty, -5)$. Numerical analysis confirms three zeros of $f$ in $(0,\pi)$. So $N_{\mathrm{int}} = 3$, $N_{\mathrm{ext}} = 1$, and $N_{\mathrm{real}} = 4$. Indeed, the four real roots are $t = -1, -2, -3, 6$.
\end{example}

\begin{example}[Four complex roots]\label{ex:four_complex}
Consider $P(t) = t^4 - 2t^2 + 3 = 0$. Here $m = -2$, $p = 0$, $q = 3$, so $u = \sqrt{2}$, $a = 0$, and $b = 8 \cdot 3/4 - 1 = 5$. Since $b = 5 > |a| + 1 = 1$, Theorem~\ref{thm:four_complex} immediately gives: all four roots are complex. Indeed, substituting $s = t^2$ gives $s^2 - 2s + 3 = 0$, so $s = 1 \pm i\sqrt{2}$, confirming that all roots are non-real.
\end{example}

\begin{example}[Two real, two complex roots]\label{ex:two_real}
Consider $P(t) = t^4 - 4t^2 + t + 1 = 0$. Here $m = -4$, $p = 1$, $q = 1$, so $u = 2$. The parameters are $a = 8/8 = 1$ and $b = 8/16 - 1 = -0.5$. The boundary values are $f(0) = 1.5 > 0$ and $f(\pi) = -0.5 < 0$. Since $f(\pi) < 0$, there is one exterior root in $(-\infty, -2)$; numerical analysis gives exactly one interior zero, so $N_{\mathrm{real}} = 2$. The quartic has roots $t \approx -2.115, 1.544, 0.351 \pm 0.710i$.
\end{example}

\section{The biquadratic case}

When $p = 0$, the quartic becomes the biquadratic $t^4 + mt^2 + q = 0$, which factors as a quadratic in $t^2$. In our framework, $a = 0$ and $b = 8q/m^2 - 1$, so $f(\theta) = \cos 4\theta + b$. The zeros in $[0,\pi]$ satisfy $\cos 4\theta = -b$, which has solutions if and only if $|b| \leq 1$, i.e., $0 \leq q \leq m^2/4$. When this holds, there are four zeros (four real roots) if $0 < q < m^2/4$; when $q = 0$, two zeros merge (double root at $t = 0$); when $q = m^2/4$, zeros merge pairwise. This matches the direct analysis via the quadratic formula in $t^2$.

\section{The case $m \geq 0$}\label{sec:mgeq0}

When $m \geq 0$, the substitution $t = \sqrt{-m}\cos\theta$ is not available. However, a hyperbolic substitution $t = \sqrt{m}\cosh\phi$ could in principle match the coefficients via the identity $8\cosh^4\!\phi - 8\cosh^2\!\phi + 1 = \cosh 4\phi$, though the unbounded range of $\cosh$ makes the analysis less clean. Alternatively, one may observe that when $m > 0$, the function $P(t)$ has $P''(t) = 12t^2 + 2m > 0$ for all $t$, so $P$ is \emph{globally} convex and therefore has at most two real roots. The nature of the roots in this case is efficiently determined by Sturm's theorem or the classical discriminant.

\section{Discussion}

\subsection{Comparison to the discriminant}

The classical quartic discriminant is a polynomial of degree six in the coefficients $m, p, q$ of the depressed quartic~\cite{Rees1922}. Computing and interpreting it is a substantial exercise. Our trigonometric method replaces this with two elementary computations: the parameters $a$ and $b$ from (\ref{eq:ab}), and the signs of $f$ at the boundary points and critical points. The trade-off is that our method requires $m < 0$, while the discriminant works for all coefficient values.

\subsection{Extension to the quintic}

The Chebyshev identity $T_n(\cos\theta) = \cos n\theta$ exists for every $n$, so the same idea can in principle be applied to polynomials of any degree. For the quintic, the identity $16\cos^5\!\theta - 20\cos^3\!\theta + 5\cos\theta = \cos 5\theta$ leads to a reduced equation involving $\cos^2\!\theta$, $\cos\theta$, and $\cos 5\theta$, whose analysis is only slightly more involved.

\subsection{Why not centuries ago?}

The Chebyshev identity (\ref{eq:cheb4}) has been known since the 18th century, and the depressed quartic has been studied since the 16th century. The substitution $t = u\cos\theta$ is a standard device in the theory of cubic equations (where it produces Vieta's trigonometric solution). Why, then, wasn't this connection to the quartic observed long ago?

Perhaps the reason is that the quartic's algebraic solution via resolvents was already known (due to Ferrari, Descartes, and Euler), and the discriminant-based classification was considered settled. The trigonometric approach provides no new \emph{computational} information---it does not solve the quartic---but it does provide a new \emph{conceptual} lens through which the root structure becomes geometrically transparent. In the spirit of~\cite{Loh2019}, where a simple proof of the quadratic formula was rediscovered by returning to basics, we hope that this trigonometric perspective may encourage the reader to look at familiar objects with fresh eyes.

\begin{thebibliography}{99}

\bibitem{Irving2013}
R.~S.~Irving,
\emph{Beyond the Quadratic Formula},
Mathematical Association of America, Washington, DC, 2013.

\bibitem{Loh2019}
P.-S.~Loh,
A simple proof of the quadratic formula,
\emph{arXiv:1910.06709}, 2019.

\bibitem{Nickalls2009}
R.~W.~D.~Nickalls,
The quartic equation: invariants and Euler's solution revealed,
\emph{Math.\ Gazette} \textbf{93} (2009), 66--75.

\bibitem{Rees1922}
P.~K.~Rees,
The discriminant of a quartic equation,
\emph{Amer.\ Math.\ Monthly} \textbf{29} (1922), no.~4, 150--153.

\bibitem{Viete1646}
F.~Vi\`ete,
\emph{Opera Mathematica},
Leiden, 1646.
(For a modern treatment of Vieta's formulas, see D.~Cox, J.~Little, and D.~O'Shea,
\emph{Ideals, Varieties, and Algorithms}, 4th ed., Springer, 2015.)

\end{thebibliography}

\end{document}
